\section{Experiments}
\label{sec:experiments}
We organize evaluation around three questions: whether query identity must be
resolved before fusion, whether repairing reference scope improves complete
extent, and whether better reconstruction or a different readout alone explains
the gains. Comparisons establish task performance; matched interventions address
these mechanism questions.

\subsection{Protocol, inputs and supervision}
The established method covers 31 scene jobs and five task outputs. Four LERF
scenes supply two outputs, so these are not five independent datasets.
All values are percentages. The benchmark rooms have been used during
development; neither prediction-before-scoring seals nor scene-level intervals
make this an untouched final test.
\begin{table}[ht]
\centering\small
\begin{tabular}{lp{.25\linewidth}p{.43\linewidth}}
\toprule
Task & Cohort and primary metric & Input and evaluation contract \\
\midrule
LERF-2D & 4 scenes, 208 observations; mIoU and localization & Source-assisted text; semantic response rendered before thresholding. \\
LERF-3D selection & Same scenes; projection mIoU and Acc.@0.25 & Independent selected Gaussians projected to 2D GT, not volumetric GT IoU. \\
ScanNet OVS & Paper8 Gaussian mIoU/mAcc; internal10 mesh mIoU & Fixed 19/15/10 menus; Gaussian and mesh domains kept separate. \\
NVOS & 8 tasks, 7 scene clusters; IoU and AP & Original reference/source access and target-image extensions explicitly separated. \\
UQIS & 9 rooms, 67 instances, 268 queries; set IoU and AP & Text, image, 2D point and 3D point; map-only query-time access. \\
\bottomrule
\end{tabular}
\caption{Evaluation units. LERF and UQIS average within scenes then across
scenes; NVOS averages tasks. Its LUDVIG reproduction first averages three seeds
within each task.}
\label{tab:protocols}
\end{table}

The implementation recorded as v10 is our method, including native matching,
admission, fusion and complete membership. LERF uses SAM3 source masks and
SigLIP2 matching; ScanNet retains SigLIP2 and Mask-Adapter priors with
literal-category SAM3 evidence. Original NVOS uses positive-derived SAM3
reference/source tracks. All UQIS maps retain SigLIP2 giant and DINOv2 ViT-B/14.
Its complete-member decoder uses 22 independently labelled training rooms and
three calibration rooms, disjoint from the benchmark. Frozen teachers do not
make this supervised component training-free. The reference-completion
extension adds original negative scribbles and DINOv2; its target-conditioned
variant also uses observed target RGB. It does not replace the input contract
of the original five-output results.

Table origins are \textbf{P}: published, \textbf{R}: locally reproduced,
\textbf{O}: our established implementation, and \textbf{M}: a measured extension.
Published values retain their reporting-paper provenance. Local aggregates
are recomputed from hash-bound reports; no new external training is performed.
No per-scene or per-task winners are combined. The full experiment catalog,
secondary variants, both NVOS GT conventions and optional complete-response
replacement are retained in the supplements.

\subsection{Five-output performance and protocol checks}
\paragraph{LERF semantic response and actual selection.}
Tables~\ref{tab:lerf2d}--\ref{tab:lerf3d} report semantic IoU/localization and
selected-set projection separately. Our method achieves 70.40 semantic IoU
and 66.96 projection IoU. SAM3 text-conditioned evidence and native SigLIP2
matching differ from the external methods' teacher access; these differences
cannot be assigned to the admission mechanism alone. The VALA reproduction
uses compatible Occam RGB geometry rather than retraining VALA geometry.
An independent CPU audit exactly reproduces all 208 selected-mask IoUs;
semantic records are reaggregated without claiming a fresh render.
\begin{table}[ht]
\centering\small\begin{tabular}{llrr}
\toprule
Method & Origin & mIoU & Loc. acc. \\
\midrule
LERF \cite{vala} & P & 37.40 & 73.60 \\
LangSplat \cite{vala} & P & 51.40 & 84.30 \\
GAGS \cite{vala} & P & 54.10 & 81.70 \\
LangSplatV2 \cite{vala} & P & 59.90 & 84.10 \\
Occam's LGS \cite{vala} & P & 61.30 & 82.50 \\
VALA \cite{vala} & P & 61.70 & 86.40 \\
Occam's LGS & R & 63.62 & 82.85 \\
Ours (v10) & O & 70.40 & 92.42 \\
\bottomrule
\end{tabular}

\caption{LERF-2D. Published context from VALA Table 1~\cite{vala}.
Localization is peak-in-box accuracy, not ranking AP. Teacher access differs.}
\label{tab:lerf2d}
\end{table}
\begin{table}[ht]
\centering\small\begin{tabular}{llrr}
\toprule
Method & Origin & mIoU & Acc.@0.25 \\
\midrule
OpenGaussian \cite{vala} & P & 38.36 & 51.43 \\
Dr. Splat \cite{vala} & P & 43.29 & 64.30 \\
InstanceGaussian \cite{vala} & P & 43.87 & 61.09 \\
CAGS \cite{vala} & P & 50.79 & 69.62 \\
VoteSplat \cite{vala} & P & 50.10 & 67.38 \\
Occam's LGS \cite{vala} & P & 47.22 & 74.84 \\
VALA \cite{vala} & P & 58.02 & 82.85 \\
VALA (compatible geometry) & R & 54.12 & 79.35 \\
Ours (v10) & O & 66.96 & 90.02 \\
\bottomrule
\end{tabular}

\caption{LERF selected-Gaussian projection, using 2D masks as GT.
Acc.@0.25 counts projections with IoU above 0.25. Published values from
VALA Table 1~\cite{vala}; its local reproduction uses compatible RGB geometry.}
\label{tab:lerf3d}
\end{table}

\paragraph{ScanNet open-vocabulary segmentation.}
Table~\ref{tab:scannet} separates eight-scene Gaussian scoring from ten-scene
mesh scoring. Our unchanged labels achieve 66.95 mIoU19 under the paper-eight
Gaussian metric and 61.01 under ten-room mesh scoring. The former applies the
VALA evaluator with anisotropic Mahalanobis pseudo labels, five-scale radius,
top-1000 class-balanced votes and opacity--volume weights. Predictions are
sealed before target-label scoring. Independent scoring reproduces all ten
mesh scenes and eight Gaussian scenes; all 1,134,207 pseudo labels agree,
including 256 checks with the released VALA implementation. On the \emph{same}
eight rooms, mesh and Gaussian mIoU19 are 61.0567 and 66.9480: that 5.8913-point
change is caused by the metric, not a learned improvement. The metric bridge
still uses provided unlabelled mesh geometry, whereas VALA uses RGB-optimized
Gaussians. Geometry and teacher access remain comparison qualifications.
\begin{table}[ht]
\centering\small\resizebox{\linewidth}{!}{\begin{tabular}{llrrrrrr}
\toprule
Method & Origin & IoU19 & Acc19 & IoU15 & Acc15 & IoU10 & Acc10 \\
\midrule
\midrule\multicolumn{8}{l}{\textit{Eight scenes: Gaussian scoring (input differences disclosed)}} \\
LangSplat \cite{vala} & P & 2.45 & 8.59 & 3.45 & 13.21 & 6.48 & 21.89 \\
OpenGaussian \cite{vala} & P & 27.73 & 42.01 & 29.67 & 46.15 & 39.93 & 57.34 \\
Dr. Splat \cite{vala} & P & 29.31 & 47.68 & 33.25 & 54.33 & 44.19 & 65.19 \\
Occam's LGS \cite{vala} & P & 31.93 & 48.93 & 34.25 & 53.71 & 45.16 & 64.39 \\
VALA \cite{vala} & P & 32.11 & 50.05 & 35.10 & 54.77 & 46.21 & 65.61 \\
VALA & R & 34.53 & 51.59 & 37.96 & 56.77 & 47.36 & 67.47 \\
Ours (v10, metric bridge) & O & 66.95 & 77.94 & 67.86 & 79.61 & 78.07 & 87.89 \\
\midrule\multicolumn{8}{l}{\textit{Ten scenes: mesh scoring (separate evaluation track)}} \\
Ours (v10) (10-mesh) & O & 61.01 & -- & 62.53 & -- & 72.49 & -- \\
Ours + replacement (L11) (10-mesh) & M & 61.01 & -- & 62.53 & -- & 72.49 & -- \\
\bottomrule
\end{tabular}
}
\caption{ScanNet OVS. Published context from VALA Table 2~\cite{vala}.
Paper8 Gaussian and internal10 mesh scoring are separate tracks.
The metric bridge aligns scoring, not teacher or geometry access.}
\label{tab:scannet}
\end{table}

\paragraph{NVOS reference responses.}
Table~\ref{tab:nvos} shows the established response, both fixed reference
admission policies and the highest-IoU refinement as an explicit tradeoff.
The primary feasibility policy reaches 92.94 IoU; the source-quality secondary
reaches 93.17 IoU and 98.95 AP. Direct image evidence reaches 93.50 IoU but is
not substituted for the actual re-rendered 3D result. Original v10 uses cached
source evidence without online target RGB; these extensions use target RGB,
negative reference scribbles and DINOv2. The improvement over 86.51 therefore
is not an input-identical ablation.

The local LUDVIG-SAM reproduction reaches 91.26 over eight tasks and three
seeds. Its GT mask-maximum convention differs from our retained $g>127$
convention; changing only this convention takes original v10 from 86.5134 to
86.6835, far below the observed extension gain. Changing bilinear
implementations moves the result by less than 0.00002 points. These checks
limit scorer explanations but do not make different input protocols identical.
Recent published rows are contextual: VCAR uses seven targets, and LEGO's
seven-scene target roster has not been verified against ours. We do not claim
same-protocol SOTA.
\begin{table}[ht]
\centering\small\begin{tabular}{llrr}
\toprule
Method & Origin & IoU & AP \\
\midrule
NVOS \cite{ludvig} & P & 70.10 & -- \\
SA3D-TRF \cite{ludvig} & P & 90.30 & -- \\
SA3D-GS \cite{ludvig} & P & 92.20 & -- \\
SAGA \cite{ludvig} & P & 92.60 & -- \\
OmniSeg3D \cite{ludvig} & P & 91.70 & -- \\
LUDVIG-DINOv2 \cite{ludvig} & P & 92.40 & -- \\
LUDVIG-SAM \cite{ludvig} & P & 91.30 & -- \\
LUDVIG-SAM2 \cite{ludvig} & P & 91.30 & -- \\
VCAR \cite{vcar} & P & 93.50 & -- \\
LEGO \cite{lego} & P & 94.20 & -- \\
WildSeg3D \cite{wildseg3d} & P & 94.10 & -- \\
Trace3D \cite{trace3d} & P & 92.50 & -- \\
LUDVIG-SAM & R & 91.26 & -- \\
Ours (v10) & O & 86.51 & 98.02 \\
Ours + reference admission (primary) & M & 92.94 & 98.86 \\
Ours + reference admission (secondary) & M & 93.17 & 98.95 \\
Ours + local decoding (tradeoff) & M & 93.28 & 98.47 \\
\bottomrule
\end{tabular}

\caption{NVOS comparison context. P values retain LUDVIG Tables 1/D and
later papers' provenance~\cite{ludvig,vcar,lego,wildseg3d,trace3d}.
N14 policies are fixed across all eight tasks. N16 has the highest measured
native IoU but lower AP (98.47 versus 98.95); it is not a uniform improvement.
All intervening variants remain in the diagnostic supplement.}
\label{tab:nvos}
\end{table}

\paragraph{Multimodal UQIS.}
Table~\ref{tab:uqis} evaluates the local ScanNet-UQIS9 v0.3.1 protocol. None
of the cited external tables evaluates these same 268 queries, so we do not
fabricate an external row. Regional readout provides an internal capability
reference; the paired component experiments below provide finer attribution.
An optional complete-response selector improves image IoU by 1.088 points,
with one query improved and other outputs unchanged. Its controls and limited
transfer are documented separately, rather than used to explain v10's gains.
\begin{table}[ht]
\centering\small\resizebox{\linewidth}{!}{\begin{tabular}{llrrrrrrrr}
\toprule
Method & Origin & T-IoU & T-AP & I-IoU & I-AP & P2-IoU & P2-AP & P3-IoU & P3-AP \\
\midrule
Ours (regional readout) & O & 12.11 & 23.69 & 26.55 & 37.75 & 42.09 & 52.74 & 41.52 & 52.36 \\
Ours (v10) & O & 16.39 & 26.60 & 39.90 & 59.31 & 49.74 & 75.79 & 49.05 & 75.33 \\
\bottomrule
\end{tabular}
}
\caption{All four UQIS modalities: text (T), image (I), 2D point (P2) and
3D point (P3). Actual-set IoU and continuous-response AP are distinct.
Both rows use our system; the regional baseline is not an external method.}
\label{tab:uqis}
\end{table}

\subsection{Q1: does query identity need to precede response formation?}
\begin{table}[ht]
\centering\small\resizebox{\linewidth}{!}{\begin{tabular}{lrrr}
\toprule
Controlled intervention & Before & After & $\Delta$ (pp) \\
\midrule
Query-conditioned identity admission & 64.12 & 66.66 & +2.54 \\
Admitted source response fusion & 50.59 & 70.38 & +19.79 \\
Native local correspondence & 30.74 & 34.88 & +4.14 \\
Complete-carrier membership reuse & 34.88 & 39.90 & +5.02 \\
\bottomrule
\end{tabular}
}
\caption{Historical matched component controls in our method. The first two
rows measure LERF selected projection and semantic IoU, respectively; the last
two measure UQIS actual image-instance IoU. These interventions are non-additive
and are not final-system leave-one-module-out experiments.}
\label{tab:core-components}
\end{table}
Query-relative identity admission improves projected selection by 2.54 points
with geometry, native features and the set solver fixed. Using admitted source
responses improves semantic IoU by 19.79 points with the evidence, map and
independent selected set fixed. This isolates whether available evidence is
used in the answer, not the novelty of a weighted mean.
Native local image correspondence adds 4.14 points with the same teachers and
extent head; complete-carrier membership adds a further 5.02 in its paired
identity-fixed contrast. The two controls distinguish matching the requested
instance from retrieving all its members.

Additional controls give a 2.35-point category gain from dual-native priors on
three ScanNet rooms, and 9.62-point gains for each UQIS point modality from
real-observation membership training. These include teacher capability and
training-observation changes, respectively. They are not evidence for a
teacher-independent weighting formula; details remain in the supplement.

\subsection{Q2: does completing reference scope before transport recover extent?}
\begin{table}[ht]
\centering\small\begin{tabular}{llrrrr}
\toprule
Matched pipeline & Orig. IoU & Compl. IoU & $\Delta$ & Orig. AP & Compl. AP \\
\midrule
Two-frame transport (N12) & 88.20 & 91.48 & 3.28 & 98.35 & 98.78 \\
Source-only fusion (N15) & 88.01 & 92.75 & 4.74 & 98.26 & 98.81 \\
\bottomrule
\end{tabular}

\caption{Reference-scope interventions on all eight NVOS tasks, retained GT.
Each paired row keeps its downstream propagation, fusion, geometry and readout
fixed. Completion introduces signed reference evidence and DINOv2 into response
construction; it is a mechanism bundle, not an isolated teacher ablation.
Source-only fitting excludes target-image evidence.}
\label{tab:reference-scope}
\end{table}
Under the same frozen two-frame memory and native-ray fitting, replacing the
original reference with the constrained completed reference raises IoU from
88.20 to 91.48 (N12). Semantic completion without the signed constraint yields
88.52; unrestricted semantic similarity does not reliably preserve instance
identity. Reference-only feasibility and quality admission subsequently reach
92.94 and 93.17 under their fixed policies (N14). Their candidates and selection
criteria are independent of target GT; design nevertheless follows feedback
on this opened cohort.

More directly, N15 propagates original and completed references through the
same 318 source views, excludes target-image evidence and applies the same
fusion. IoU rises from 88.01 to 92.75, with AP rising from 98.26 to 98.81.
Five tasks improve and three regress; completion is not uniformly beneficial.
This paired result shows that the scope intervention survives source-view
fusion and is not explained solely by decoding the target image. It does not
separately identify the contribution of negative scribbles, DINOv2 and frozen
SAM3 completion. The central supported claim is about the completed-reference
mechanism under this declared access contract.

\subsection{Q3: are stronger fitting or additional decoding sufficient?}
With fixed evidence, squared reconstruction reduces its source objective on
all 22 source-assisted jobs. Yet LERF semantic IoU/AP change from 70.40/87.62
to 70.80/87.00, and ScanNet mIoU19 falls from 61.01 to 60.28. UQIS outputs are
unchanged compatibility references. Better fit to admitted labels therefore
does not establish better task responses across outputs.

Within the completed-reference experiment, adding target evidence to the source
fit gives 93.03 IoU and 98.98 AP. Against the target-only N14 response,
source residual decreases from 0.004626 to 0.000748 while IoU drops by 0.14
points and AP gains 0.03. All tasks reduce both true and false positives:
consistency suppresses errors but also useful extent. This is a tradeoff,
not evidence that a lower residual solves scope correctness.

Query-local target re-decoding gives 93.28 IoU, only 0.1075 above its N14
control, while AP loses 0.4787; four tasks improve and four regress.
Removing the horns scene cluster eliminates the mean IoU advantage.
A numerical audit exactly recovers all 24 image arrays and checks stable
probability-mixture ranking: the largest macro AP correction is below 0.002
points and no hard masks change. Thus the image AP loss is not explained by
sigmoid saturation or interpolation precision. This audit does not rerun the
3D fits. The joint-fusion and re-decoding search routes are stopped rather than
promoted on one metric.

Finally, retained field-by-readout controls report original NVOS semantic IoU
86.51 and independent selected-set projection IoU 85.22. Crossed outputs are
scored separately. The analysis in Section~\ref{sec:analysis} explains why one
cannot replace the semantic result with a selected-set number or infer better
actual selection from a modified score field. The supplements retain all
crossed readouts, solver failures, prompt controls and per-task diagnostics.
